replacement. An exceptional side can occur for only one step because its
$A$-successor is ordinary. Thus an equal-rank passage strictly decreases
$\tau$ and must reach the exponent-one/two factor boundary. No new witness is
introduced during this countdown.
\end{proof}

\begin{lemma}[Factor and marked-tail routing]\label{lem:factor-routing}
Every factor-fork, high-return, marked-tail, short-lift, and terminal control
has one of the following effects:
\begin{enumerate}[label=(\alph*)]
\item strict inner-source decrease;
\item strict replacement of an existing inner token by certified lower
      descendant token(s);
\item a control change listed in Table~\ref{tab:control} with the retained
      source and token multiset unchanged; or
\item a same-control recurrence with strictly smaller local counter $\tau$.
\end{enumerate}
\end{lemma}

\begin{proof}
The boundary and loss-sibling controls are covered by Lemmas~\ref{lem:boundary-normalize} and~\ref{lem:loss-sibling}.

The factor exponent-one row uses the shared exponent-one/two child. If the
selected finite source is ordinary, Lemma~\ref{lem:side} exposes a common
child of the retained \LOSS\ witness and hence a lower \WIN\ token; the
exceptional row is an exact constant-tail lift over the same \LOSS\ witness.
For factor exponent at least two, the reverse constant-tail parent either
removes a factor at exponent at least two or exposes a finite parent-child
witness; in the latter case its common child is a certified descendant.

A high return has valuation $5,6$, or at least $7$. In the first two rows the
explicit short-tail formulas replace the carried return token before the new
marked task is installed. For valuation at least $7$, the universal
three/two frame carries a \WIN/LOSS\ pair; the selected common child is a
certified descendant before the tail is restarted. Hence every high-return
edge is strict in $\Psi(N)$.

For a marked constant tail of length $D\ge3$, if its marked states are
$q$ (\WIN) and $y$ (\LOSS), the common child
\[
  r=B(q)=A(y)
\]
for $D\ge4$ (and the boundary common child for $D=3$) is a \WIN\ child of
$y$, so $h(r)<h(y)$. Thus the existing token $y$ is replaced before any new
inner task is installed. The two short rows $D=1,2$ are exact lifts over an
already carried token and merely change control. Long exact-lift recurrences
strictly lower their remaining tail/factor exponent, which is $\tau$.
Terminal exponent-two/three macros reconstruct a finite descendant token
before they restart any obligation. The strict token replacements used in these rows are collected explicitly in Table~\ref{tab:tokens}. These cases exhaust the factor/tail constructors because their only unbounded parameters are the valuation and tail intervals just split above.
\end{proof}

The equal-retained-rank control edges and a concrete control height $c(v)$ are
shown in Table~\ref{tab:control}. Any edge not listed there is strict in an
earlier component of \eqref{eq:full-rank}. Same-control tail/factor recurrences
strictly decrease $\tau$.

\begin{table}[ht]
\centering
\small
\begin{tabular}{@{}lll@{}}
\toprule
Control $v$ & $c(v)$ & Equal-rank successors \\
\midrule
\texttt{marked\_tail} & 9 & \texttt{short\_lift} \\
\texttt{boundary\_entry} & 8 & \texttt{a\_obligation}, \texttt{factor\_fork} \\
\texttt{short\_lift} & 8 & \texttt{a\_obligation} \\
\texttt{a\_obligation} & 7 & \texttt{a\_test\_4}, \texttt{factor\_fork} \\
\texttt{a\_test\_4} & 6 & \texttt{a\_test\_3}, \texttt{b\_select} \\
\texttt{a\_test\_3} & 5 & \texttt{a\_test\_2}, \texttt{b\_select} \\
\texttt{a\_test\_2} & 4 & \texttt{a\_test\_1}, \texttt{b\_select} \\
\texttt{a\_test\_1} & 3 & \texttt{b\_select} \\
\texttt{b\_select} & 2 & \texttt{b2\_first} \\
\texttt{loss\_sibling\_entry} & 2 & \texttt{factor\_fork} \\
\texttt{b2\_first} & 1 & \texttt{b2\_ready} \\
\texttt{factor\_fork} & 1 & \texttt{high\_return} \\
\texttt{b2\_ready} & 0 & -- \\
\texttt{high\_return} & 0 & -- \\
\texttt{terminal\_macro} & 0 & -- \\
\bottomrule
\end{tabular}
\caption{Equal-rank control graph. Every listed edge strictly lowers $c(v)$.
Strict source/token/bit edges may reset $c$ and $\tau$.}
\label{tab:control}
\end{table}

\begin{proposition}[Well-founded typed normalizer]\label{prop:typed-normalizer}
The typed routing relation is well-founded for every typed entry.
\end{proposition}

\begin{proof}
Consider the rank \eqref{eq:full-rank}. A retained-source decrease is strict in
the first component and may reset all later data. At fixed retained source,
$\eta:1\to0$ is strict and may seed $M$ and all later routing data. Once
$\eta=0$, every change of $M$ is a certified strict descendant replacement
and therefore decreases $\Phi(M)$ by Lemma~\ref{lem:token-rank}; it may reset
all later inner data.

At fixed earlier components, $\zeta:1\to0$ is a one-time strict inner-token
seed. Thereafter every change of $N$ strictly decreases $\Psi(N)$. With the
token components fixed, a promoted inner-source comparison strictly lowers
$a$. If all preceding components are unchanged, the semantic routing lemmas
above show that a control-changing equal-rank edge is one of the edges of
Table~\ref{tab:control}, hence strictly lowers $c(v)$; a same-control
tail/factor recurrence strictly lowers $\tau$. Thus every transition
strictly decreases one component while preserving all earlier components.
The lexicographic product of these well-orders is well-founded.
\end{proof}

\section{Proof of the main theorem}

\begin{proof}[Proof of Theorem~\ref{thm:main}]
Assume that a \DRAW\ exists and choose the least state source $s$ occurring
among \DRAW\ positions. Begin with $\eta=1$; no ranked token has yet been installed.

If $s=0$, Lemma~\ref{lem:zero-source} performs a finite pre-entry
normalization and reaches a boundary or loss-sibling typed entry. Consume
$\eta:1\to0$ exactly there and invoke Proposition~\ref{prop:typed-normalizer}.

Let $s>0$. Follow an actual \DRAW\ child through long constant tails until
the tail exponent is one or two. Proposition~\ref{prop:boundary-reduction}
then gives a typed entry, a strict retained-source exit, or the factor-free
exponent-two state $Q_2^e(J(s))$ \DRAW. The last case is handled by
Proposition~\ref{prop:base-entry}. A strict smaller-source exit contradicts
the choice of $s$; otherwise a typed entry is obtained and $\eta$ is
consumed. Thus all four boundary combinations $(r,k)=(1,0),(1,>0),(2,0)$ and
$(2,>0)$ are covered.

From then on every macro transition lies in the well-founded relation of
Proposition~\ref{prop:typed-normalizer}, unless an earlier outer component
strictly decreases, in which case the lexicographic rank decreases even more
strongly. Yet every \DRAW\ has no \LOSS\ child and at least one \DRAW\ child.
Therefore from every marked \DRAW\ configuration an outcome-compatible
continuing \DRAW\ can be chosen, yielding another finite macro transition.
Iteration would create an infinite descending chain in the rank
\eqref{eq:full-rank}, impossible. Hence no \DRAW\ exists in the conjugated
game.

It remains only to transfer the conclusion back to an arbitrary original
odd starting state. Write it as $n=2m+1$. If $m=0$, then $n=1$ is terminal.
For $m>0$, Proposition~\ref{prop:first-normal} says that its two original
children are $F(m)$ and $F(R(m))$. These are represented in the conjugated game
by the states $m$ and $R(m)$, respectively, and both have finite remoteness by
the no-\DRAW\ conclusion just proved. Hence the original state $n$ is itself
\WIN\ or \LOSS\ with finite height. Thus every positive odd starting position
has finite remoteness and optimal play reaches $1$.
\end{proof}

\section{Verification boundary and reproducibility}

The theorem above is a human mathematical claim. The accompanying files have
narrower roles.
\begin{itemize}[leftmargin=*]
\item \texttt{verify\_v3\_0.py} checks the long-tail identities, fixed-tail
      diamonds, ordinary side relation, base-entry valuations and inequalities,
      source-zero valuation table, hidden-parent identity, and the finite
      equal-rank control graph over large finite ranges.
\item \texttt{routing\_certificate\_v3\_0.json} records the control states,
      equal-rank edges, strict edge classes, and rank components used in
      Table~\ref{tab:control}. It is a structural certificate, not a substitute
      for the semantic proofs in this article.
\item The repository's Lean development checks selected definitions and
      general well-founded metatheory. It is not claimed to kernel-check the
      complete concrete theorem.
\end{itemize}

The finite computations support the arithmetic and the transcription of the
routing table; the infinite conclusion rests on the proved identities,
outcome recursion, and the well-founded rank.

\appendix
\section{Exceptional side formulas}\label{app:exceptional}

For completeness, the four exceptional classes of Lemma~\ref{lem:side} have
exact formulas. For $t\ge0$,
\[
\begin{array}{c|c|c}
q&B(q)&B(A(q))\\ \hline
16t+1  &R(6t)   &18t+1\\
16t+3  &R(6t+1) &18t+4\\
16t+12 &R(6t+4) &18t+13\\
16t+14 &R(6t+5) &18t+16.
\end{array}
\]
The displayed $B(A(q))$ value always has constant-tail length one. Hence an
exceptional side step cannot repeat immediately; its next $A$-source is
ordinary. These formulas also justify the exceptional branch used in
Proposition~\ref{prop:base-entry}.

\section{Routing inventory}\label{app:routing}

The equal-rank graph in Table~\ref{tab:control} is deliberately smaller than
the full outcome case split: source and token exits are omitted from the graph
because they are already strict in earlier components. The semantic partitions
are as follows.
\begin{center}
\small
\begin{tabular}{@{}p{.28\linewidth}p{.62\linewidth}@{}}
\toprule
Control & Exhaustive nonterminal partition \\
\midrule
\texttt{boundary\_entry} & source/token exit,
  $\texttt{a\_obligation}$, or $\texttt{factor\_fork}$ after finite tail countdown \\
\texttt{loss\_sibling\_entry} & source/token exit, or
  $\texttt{factor\_fork}$ after finite tail countdown \\
\texttt{a\_obligation} & source exit, token exit, canonical $A$ test, or
  retained-witness factor fork \\
\texttt{a\_test\_4,3,2,1} & source exit, token exit, next test, or
  $B$-selecting switch; the last test has no next-test row \\
\texttt{b\_select} & valuation $2$, or valuation at least $3$ (strict source) \\
\texttt{b2\_first} & strict source, or one-time witness installation \\
\texttt{b2\_ready} & strict source, or strict witness descendant \\
\texttt{factor\_fork} & exponent-one source/token, higher source/token, or
  adjacent high return \\
\texttt{high\_return} & valuations $5,6$ or at least $7$; each spends a token
  before marked-tail restart \\
\texttt{marked\_tail} & $D=1$, $D=2$, or $D\ge3$; long row replaces the
  marked \LOSS\ token by its common \WIN\ child \\
\texttt{short\_lift} & exact lift into an obligation \\
\texttt{terminal\_macro} & exponent-two/three terminal rows; each begins
  after a strict token replacement \\
\bottomrule
\end{tabular}
\end{center}
This inventory is exactly the case structure used in the proof of
Proposition~\ref{prop:typed-normalizer}; there is no implicit ``other'' row.

\begin{table}[ht]
\centering
\small
\begin{tabular}{@{}p{.23\linewidth}p{.29\linewidth}p{.37\linewidth}@{}}
\toprule
Routing row & Retained replacement & Certified comparison \\
\midrule
ordinary factor fork & $b\mapsto X$ & $X$ is a child of the canonical
  \LOSS\ child of $b$, hence $h(X)\le h(b)-2$ \\
exceptional factor fork & $b\mapsto \ell$ & $\ell$ is the actual \LOSS\
  child of $b$, hence $h(\ell)=h(b)-1$ \\
high return & $(u,c)\mapsto(p,b)\mapsto(q,y)$ & each new state is an actual
  descendant of the token it replaces; each replacement is strict \\
marked tail $D\ge3$ & $y\mapsto r$ & $r$ is a child of the retained
  \LOSS\ state $y$, so $h(r)\le h(y)-1$ \\
terminal exponent $2/3$ & $\ell\mapsto Z_2$ or $Z_3$ & the reconstructed
  terminal token is a proper proof-tree descendant of $\ell$ \\
zero-recycle exponent $3$ & $y\mapsto r$ & $r$ is a child of the retained
  \LOSS\ token $y$ \\
\bottomrule
\end{tabular}
\caption{Token replacements used by the strict routing rows. A branching
exceptional diamond may replace one token by several descendants; in that
case Lemma~\ref{lem:token-rank} applies to the resulting multiset.}
\label{tab:tokens}
\end{table}

\begin{thebibliography}{9}
\bibitem{Guy1986}
R. K. Guy, John Isbell's Game of Beanstalk and John Conway's Game of
Beans-Don't-Talk, \emph{Mathematics Magazine} 59(5) (1986), 259--269.

\bibitem{Guy1996}
R. K. Guy, Unsolved Problems in Combinatorial Games, Problem 42, in
\emph{Games of No Chance}, MSRI Publications 29, Cambridge University Press,
1996.

\bibitem{Althofer2024}
I. Alth\"ofer, M. Hartisch, T. Zipproth, Analysis of a Collatz Game and Other
Variants of the $3n+1$ Problem, \emph{Advances in Computer Games} (2024),
123--132.

\bibitem{AlthoferPage}
I. Alth\"ofer, Collatz Problems: The Two-Player $3n+-1$ Game,
\url{https://althofer.de/collatz-prizes.html}.

\bibitem{Repo}
G. Pochuev, Optimal $3n\pm1$ Game: Proof and Verification Repository,
\url{https://github.com/Grisha-Pochuev/3n-plus-minus-1-game}.
\end{thebibliography}

\end{document}
